El sector de la distribución de productos profesionales de peluquería se caracteriza por una fuerte dependencia de procesos manuales y herramientas de gestión poco integradas. En muchos casos, los distribuidores gestionan su actividad diaria utilizando múltiples aplicaciones independientes, hojas de cálculo o incluso procesos manuales, lo que dificulta la centralización de la información y el seguimiento eficiente de la actividad comercial.

En el caso analizado en este proyecto, la mayor parte de las tareas relacionadas con la gestión de pedidos, el seguimiento de clientes, la organización de visitas comerciales o el control de cobros se realizan mediante herramientas dispersas como hojas de cálculo, bases de datos locales y software de facturación independiente. Esta situación provoca una elevada carga administrativa y limita la capacidad de automatizar o escalar la actividad comercial.

Ante este contexto, surge la necesidad de desarrollar una solución tecnológica que permita integrar estas funcionalidades en una única plataforma digital. Una aplicación de estas características puede facilitar la gestión del negocio, reducir el tiempo dedicado a tareas administrativas y mejorar la relación entre el distribuidor y las peluquerías profesionales con las que trabaja.

Desde el punto de vista académico, este proyecto permite aplicar de forma práctica distintos conceptos estudiados durante el grado en Ingeniería Informática, como el análisis de requisitos, el modelado de sistemas, el diseño de aplicaciones web y el desarrollo de arquitecturas full-stack. Además, el desarrollo de un sistema basado en tecnologías web modernas permite explorar enfoques actuales de desarrollo de software, como las aplicaciones web progresivas (PWA) y la integración de distintos servicios en una misma plataforma.

Por todo ello, el desarrollo de este proyecto resulta relevante tanto desde el punto de vista práctico, al proponer una solución a una problemática real, como desde el punto de vista académico, al permitir aplicar y consolidar los conocimientos adquiridos durante la formación universitaria.